\begin{trapbox}
\textbf{\textcolor{CTrap}{易错 1｜题目条件本身无解，却继续套公式}}

\textbf{错法：}在 $a=5,b=4$ 的椭圆上给出 $d=2$，由 $\sqrt{r_1r_2}=ab/d$ 得到 $10$，当成有效答案。
\textbf{纠正：}先检查 $\boxed{b\le d\le a}$。这里必须有 $4\le d\le5$，故不存在该切线；且 $\sqrt{r_1r_2}\le a=5$ 也能反证数值不可行。

\medskip
\textbf{\textcolor{CTrap}{易错 2｜把“焦点距切线”当“中心距切线”}}

\textbf{错法：}把 $h_1$、$h_2$ 误作 $d$，写成 $r_1r_2h_1h_2=a^2b^2$。
\textbf{纠正：}$d=\operatorname{dist}(O,l)$，$h_i=\operatorname{dist}(F_i,l)$；正确关系是
\[
\boxed{h_1h_2=b^2},\qquad
\boxed{h_i=\frac{r_i}{a}d},\qquad
\boxed{r_1r_2d^2=a^2b^2}.
\]

\medskip
\textbf{\textcolor{CTrap}{易错 3｜把几何平均误记成算术平均}}

\textbf{错法：}写成 $\dfrac{r_1+r_2}{2}d=ab$。因为 $r_1+r_2=2a$，这会错误推出 $d=b$ 恒成立。
\textbf{纠正：}$r_1+r_2=2a$ 是定义，而 C020 中是 $\sqrt{r_1r_2}$。

\medskip
\textbf{\textcolor{CTrap}{易错 4｜焦半径一次式的正负颠倒}}

在 $F_1(-c,0),F_2(c,0)$ 的记号下，正确公式是
\[
\boxed{r_1=a+\frac{cx_0}{a}},\qquad
\boxed{r_2=a-\frac{cx_0}{a}}.
\]
当 $x_0>0$ 时，$P$ 更靠近右焦点，故 $r_2$ 更小；也可用 $P=(a,0)$ 立即核对 $r_1=a+c,r_2=a-c$。

\medskip
\textbf{\textcolor{CTrap}{易错 5｜点到直线距离公式代错点}}

$P$ 处切线为 $\frac{x_0}{a^2}x+\frac{y_0}{b^2}y-1=0$。
计算的是\textbf{原点}到切线的距离，故分子为 $|0+0-1|=1$，
\[
d=\frac1{\sqrt{x_0^2/a^4+y_0^2/b^4}}>0.
\]
若把 $P$ 代进去，结果为 $0$，那算的是切点到自身切线的距离。

\medskip
\textbf{\textcolor{CTrap}{易错 6｜光学角平分方向搞反}}

椭圆\textbf{法线}平分 $\angle F_1PF_2$ 的内角；\textbf{切线}平分与之相邻的外角。对焦点到切线距离用
\[
h_i=r_i\cos\frac\theta2,
\]
而不是 $r_i\sin(\theta/2)$。可代入最右顶点检验：此时 $\theta=0$、$h_i=r_i$，必须是 $\cos0=1$。

\medskip
\textbf{\textcolor{CTrap}{易错 7｜证明只验证两个端点}}

在 $(a,0)$、$(0,b)$ 验算乘积都是 $ab$，只能帮助\textbf{发现规律}，不能证明任意动点成立。应补几何法的余弦定理，或坐标法的变量消元。

\medskip
\textbf{\textcolor{CTrap}{易错 8｜公式适用范围误解}}

\begin{itemize}[leftmargin=*,itemsep=0.3em]
\item $a$ 始终代表\textbf{半长轴}；若长轴沿 $y$ 轴，改用 $x^2/b^2+y^2/a^2=1$ 及真实焦点即可，主公式仍然是 $ab$，无需换成别的常数。
\item 平移、旋转椭圆不改变距离与半轴，公式仍成立，但 $d$ 应从\textbf{平移后的椭圆中心}量起，不再自动是坐标原点。
\item $P$ 必须是切点，$l$ 必须是该点的切线；任意过 $P$ 的割线不满足 C020。
\item $c$ 是\textbf{半焦距}，两焦点之间的距离为 $2c$；不要把 $c$ 错当完整焦距。
\end{itemize}
\end{trapbox}
